\documentclass[../main.tex]{subfiles}
\begin{document}
%==========================================TIÊU ĐỀ TẬP========================================
\begin{table}[h]
  \begin{adjustwidth}{-1cm}{}
    \begin{tabular}{>{\centering\arraybackslash}p{6cm}|>{\raggedright\arraybackslash}p{15cm}}
      \multirow{5}{6cm}
      % Cột bên trái: ÉPISODE và số tập
      {\\[-40pt]
        \flaregothic\fontsize{20pt}{20pt}\selectfont \centering
        \color{eptype}ÉPISODE\color{black}\\
        \burbank\fontsize{72pt}{72pt}\selectfont
        \color{epnum}0\color{black}\\[6pt]
        \cafeteria\fontsize{20pt}{20pt}\selectfont\color{epnum}(0-1)\color{black}
        \\[18pt]
        \flaregothic\fontsize{14pt}{14pt}\selectfont
        \color{epattr}BIENVENUE, \uline{\color{Fubuking} Fubuki}, \uline{\color{Matsugod} Matsuri},
        \uline{\color{Robopon} Roboco}, et \uline{\color{Akutan} Aqua} !\\
        \color{black}
      }
       &
      % Dòng thứ nhất: tên tập tiếng Nhật
      {\fotskip\fontsize{18pt}{18pt}\selectfont Cooming Soon!! 【\ruby{予告}{こく}\ruby{編}{へん}】
      }                                    \\[6pt]
      % Dòng thứ hai: tên tập tiếng Anh
       & {\cafeteria\fontsize{18pt}{18pt}\selectfont The Trailer: Don't Sleep Unless You Miss It!}                           \\[4pt]
      % Dòng thứ ba: tên tập tiếng Việt
       & {\vnmsans\fontsize{18pt}{18pt}\selectfont Những cảnh quay đầu tiên}                                                 \\[8pt]
      % Dòng thứ tư: Nhân vật xuất hiện
       & {\caviar{\fontsize{14pt}{14pt}\selectfont Nhân vật: \emp{kuon1} \emp{roboco} \emp{fubuki} \emp{matsuri} \emp{aqua}}} \\[8pt]
      % Dòng cuối cùng: Thời gian diễn ra của tập (thường sẽ là ngày phát hành)
       & {\continuum\fontsize{16pt}{16pt}\selectfont Ngày phát hành: 28/04/2019}                                             \\[18pt]
    \end{tabular}
  \end{adjustwidth}
\end{table}
%=========================================TRUYỆN CHÍNH========================================
\color{black}

\txtbx[2]{
  {\color{vol} \uline{Tóm tắt:} Buổi ghi hình thử nghiệm đầu tiên tại văn phòng Cover Corp diễn ra một cách thảm họa khi ba cô nàng Aqua, Matsuri và Roboco bị một sinh vật bay bí ẩn truy đuổi khiến cả ba hoảng loạn nhảy lầu, phá hỏng cửa kính trong sự ngơ ngác của Fubuki. Toàn bộ vụ hỗn loạn được ghi lại qua ống kính một chiếc máy quay ẩn, kết thúc bằng thông điệp ``Coming Soon!!'' trên trần nhà. Sau đó, chàng tân tổng quản lý Hikawa Kuon đến nhận việc, giải cứu ba cô gái khỏi chú gián đất bằng chiếc ``dép lào huyền thoại'' và cùng họ thảo luận về dự án phim ngắn \textit{Holo no Graffiti}. Sau khi hoàn tất thủ tục ký hợp đồng cùng CEO Tanigo - YAGOO - và A-chan, Kuon chính thức đảm nhận vai trò mới, mở ra một kỷ nguyên đầy hứa hẹn nhưng cũng không kém phần quái đản tại văn phòng \textit{hololive}.}
}

\begin{center}
  \uline{(Câu chuyện được kể dựa trên một đoạn phim được ghi bằng chiếc máy quay cố định trên bàn.)}
\end{center}

Mười giờ mười phút sáng. Một ngày Chủ Nhật nọ.

Văn phòng của công ty lúc này vẫn còn chìm trong một sự ngổn ngang đầy hỗn độn. Những hộp bìa các-tông chất chồng lên nhau góc phòng, vài thùng hàng còn chưa kịp khui ra, chứng tỏ nơi này mới chỉ đi vào hoạt động cách đây vỏn vẹn vài ngày. Trên những dãy bàn làm việc, đống chai lọ rỗng, cốc nhựa và đĩa giấy nằm lăn lóc - tàn tích của một bữa tiệc mừng khai trương vô cùng linh đình vào tối hôm trước vẫn chưa được dọn dẹp. Giữa căn phòng rộng lớn và tĩnh lặng, chỉ có tiếng ``tích tắc, tích tắc'' đều đặn đầy tẻ nhạt phát ra từ chiếc đồng hồ treo tường mới cứng.

Và đây, chính là nơi mà cuộc phiêu lưu đầy những biến động bất ngờ của các cô gái \textit{hololive} cùng chàng tân Tổng quản lý Hikawa Kuon chính thức được bắt đầu.

\il

\img{0-1a}

Cánh cửa văn phòng khẽ mở ra, phá tan bầu không khí im lìm của buổi sáng Chủ Nhật. Một cô gái bước vào, mang theo năng lượng rạng rỡ của một ngày mới cùng câu chào quen thuộc: ``Chào buổi sáng mọi người\~{}!''

\model{26}{r}{6cm}{fubuki}

Thế nhưng, đáp lại tiếng gọi ngân vang đầy hào hứng của cô chỉ là một sự im lặng tuyệt đối. Căn phòng trống trơn không một bóng người.

Cô nàng ngơ ngác chớp chớp mắt, bước thêm vài bước rồi nghiêng đầu lẩm bẩm: ``Ủa? Không có ai ở đây thật sao?''

\chrint

Đó chính là Shirakami Fubuki (\ruby{白}{しら}\ruby{上}{かみ}フブキ, \ruby{Si-ra}{Bạch}-\ruby{ca-mi}{Thượng}\ Phu-bu-ki) – cô nàng Cáo trắng thuộc Thế hệ thứ Nhất kiêm trưởng nhóm GAMERS danh tiếng của \textit{hololive}. Trên sơ yếu lý lịch, Fubuki tự họa bản thân là một otaku cuồng anime và manga thứ thiệt, vô cùng năng nổ và cực kỳ thích trò chuyện với mọi người. Cô nàng thậm chí còn từng tuyên bố rằng mình đang sống trong một căn cabin bằng gỗ của ông nội, tọa lạc trên một đỉnh núi heo hút gần khu rừng huyền bí nơi loài yêu tinh sinh sống.

Bình thường năng động và hoạt ngôn là thế, nhưng sâu bên trong, Fubuki lại có một mặt tính cách khá nhút nhát mỗi khi phải đối mặt với những người mới gặp (ngoại trừ Hikawa Kuon, bằng một sợi dây liên kết kỳ lạ nào đó). Những lúc xấu hổ hay bối rối, đôi tai cáo mềm mại trên đầu cô sẽ khẽ cụp lại đầy đáng yêu. Thế nhưng, đằng sau vẻ ngoài ngọt ngào ấy lại là một ``game thủ'' thực thụ với phong cách chơi lạnh lùng, tàn nhẫn và đôi khi cực kỳ ``xoắn não'' đối thủ. Không chỉ vậy, nàng Cáo trắng này còn là một con nghiện gacha hạng nặng, sẵn sàng nướng sạch toàn bộ số tiền lương kiếm được chỉ để đổi lấy những lượt quay may rủi trong trò chơi. Fubuki cũng là kiểu người sẵn sàng đi ngược lại các quy chuẩn xã hội nếu chúng đi ngược lại với nguyên tắc và sở thích cá nhân của mình.

Fubuki đứng chững lại giữa phòng một lát, đôi tai cáo khẽ vẩy nhẹ khi cô phóng tầm mắt quét qua một lượt quang cảnh bừa bộn sau bữa tiệc hôm qua.

Nhận ra mình là người đến sớm nhất, cô khẽ thở dài, trong lòng thầm nghĩ: ``\textit{Xem ra vẫn chưa có ai tới cả nhỉ\dots\ Chắc phải một lúc nữa mọi người mới đến đông đủ. Cả Kuon-senpai nữa, hôm nay là ngày đầu tiên anh ấy nhận việc mà ta.}''

Cô nhún vai tự nhủ: ``Thôi thì cứ kiên nhẫn đợi chút vậy.''

Nghĩ đoạn, Fubuki rảo bước tiến về phía chiếc ghế sô-pha êm ái đặt ở góc phòng. Cô thả mình xuống đệm, rút điện thoại ra lướt mạng để giết thời gian, hoàn toàn không hay biết rằng, ở góc bàn đối diện, ống kính máy quay nhỏ của một ai đó vẫn đang âm thầm ghi lại từng cử chỉ nhỏ nhất của cô.

\ct

\img{0-1b}

Mười lăm phút đằng đẵng trôi qua trong im lặng. Vẫn không có thêm một bóng người nào xuất hiện. Nàng Cáo trắng vẫn ngồi một mình trên chiếc sô-pha, nhưng mí mắt cô bắt đầu trĩu nặng. Sự tĩnh lặng tuyệt đối của căn phòng như một liều thuốc ngủ tự nhiên, khiến cô không ngừng gật gù.

``Ối chà\dots\ Buồn ngủ quá đi mất\dots''

Fubuki ngáp một hơi thật dài, đặt chiếc điện thoại xuống bàn rồi vươn vai thật cao để xua đi cơn buồn ngủ đang kéo đến.

Thế nhưng, mọi nỗ lực đều vô dụng trước sự chán nản của việc phải chờ đợi một mình. Đầu cô khẽ nghiêng rồi dần ngả xuống tấm đệm sô-pha êm ái. Chỉ vài giây sau, nàng Cáo trắng đã chìm sâu vào giấc ngủ một cách ngon lành, tiếng thở khẽ đều đặn vang lên trong căn phòng vắng. Cô thầm nghĩ mọi chuyện cứ thế trôi qua yên bình cho đến khi mọi người tới\dots

Bỗng nhiên, một tiếng vo ve chói tai vang lên, xé toạc bầu không khí tĩnh lặng đến lặng người của văn phòng. Ngay sau đó, tiếng hét thất thanh đầy hoảng loạn của một cô gái từ xa vọng lại, mỗi lúc một gần. Cánh cửa phòng bị đẩy bật ra một cách thô bạo. Một cô gái lao vội vào trong như một cơn gió, gương mặt tái mét, dáng vẻ thất thần như thể đang bị truy lùng bởi một quái thú thời tiền sử.

Nhưng nếu nhìn kỹ, thứ bám đuổi ngay sau lưng cô thực chất chỉ là một sinh vật bay nhỏ bé. Bạn có thể tự hỏi: ``Thứ đó rốt cuộc là con gì vậy?'' Thế nhưng, thật trớ trêu, hình ảnh của sinh vật ấy lại bị làm mờ một cách vô cùng bí ẩn và đầy hài hước trong ống kính máy quay. Chúng ta chỉ có thể lờ mờ đoán rằng đó là một loài côn trùng có cánh – nỗi khiếp sợ kinh hoàng của đa số phái yếu, có thể là một con ruồi, một chú ong, hay một sinh vật kỳ dị nào đó mà mắt thường không thể xác định nổi.

\chrint

\model{25}{l}{7cm}{aqua}

Minato Aqua (\ruby{湊}{みなと}あくあ, \ruby{Mi-na-tô}{Tấu}\ A-cư-a) - cô nàng hầu gái tiếng tăm thuộc Thế hệ thứ Hai của \textit{hololive}, nổi tiếng với mái tóc tím hồng đặc trưng khẽ uốn xoắn tựa như một củ hành tây đáng yêu. Dù khoác trên mình bộ đồng phục hầu gái điệu đà, Aqua lại cực kỳ vụng về và thường xuyên vụt mất danh hiệu ``hầu gái chuẩn mực'' mỗi khi đụng tay vào công việc thực tế. Thế nhưng, bù lại cho sự hậu đậu ấy, cô lại là một cao thủ game đích thực với kỹ năng bắn súng góc nhìn thứ nhất (FPS) thuộc hàng thượng thừa.

Đầu óc thường xuyên bay bổng trên mây, Aqua rất dễ trở thành nạn nhân của những trò đùa tinh nghịch từ đồng nghiệp, và đôi khi là tự mình hại mình mà cô nàng chẳng hề hay biết. Biệt danh hài hước ``Baqua'' (Aqua ngố) cũng từ đó mà ra đời. Aqua còn mang tính cách khá trẻ con, dễ khóc, dễ hờn dỗi mỗi khi mọi việc không theo ý muốn, giống như lần cô nàng đùng đùng bỏ về chỉ vì bị một đàn em trêu ghẹo rằng mình không phải hầu gái thực thụ\footnote{Xem trong Quyển số Bảy, \flaregothic{ÉPISODE 94}.}

Điều trớ trêu hơn cả là dù là một streamer đình đám, Aqua lại là một người hướng nội cực kỳ nặng, vô cùng ngại giao tiếp với thế giới bên ngoài. Việc cô nàng lấy hết dũng khí để đến gặp mặt trực tiếp Kuon vào ngày hôm qua đã là một nỗ lực phi thường, mở ra một bước tiến dài trong việc giúp cô mở lòng hơn với mọi người xung quanh.

Là một nàng hầu gái thủy thủ cách tân, Aqua diện chiếc váy xếp ly màu xanh nước biển ngắn ngang đùi với phần ống tay phồng, hở nhẹ trước ngực. Đôi chân cô mang đôi giày mary janes cao gót màu xanh đậm điệu đà, đầu đội chiếc mũ hầu gái có in biểu tượng chiếc mỏ neo cá tính. Điểm xuyết trên trang phục là cổ tay áo trắng muốt đồng điệu với cổ mắt cá chân. Mái tóc màu tím hồng của cô được tết thành hai bím xoắn ốc thắt ruy-băng xanh, xen lẫn những lọn tóc với điểm nhấn màu xanh nhạt đầy thời thượng.

\chrint

\gif{0-1d}{36}{245}

Nàng `Hành tím' tội nghiệp lúc này vừa cuống cuồng chạy vòng quanh phòng, vừa khua tay múa chân loạn xạ, cái đầu lắc nguầy nguậy đầy tuyệt vọng. Cô hét lên kinh hoàng: ``\textbf{Không, không, không, không, không, không! Tránh xa ta raaaa!}''

Cứ chốc chốc, cô nàng lại nhảy dựng lên như lò xo, tìm cách xua đuổi con côn trùng đáng sợ đang lượn lờ sát sạt. Khi mọi nỗ lực xua đuổi đều đi vào ngõ cụt, trong cơn hoảng loạn tột độ, Aqua không thèm suy nghĩ mà lao thẳng người ra phía cửa sổ lớn để thoát thân. Một tiếng *RẦM* đinh tai nhức óc vang lên, tấm kính cửa sổ vỡ tan tành thành trăm mảnh thủy tinh lấp lánh rơi rào rào xuống đất. Nàng hầu gái cùng sinh vật bay đáng ghét kia lập tức biến mất sau khung cửa sổ trống hoác.

Tiếng động lớn chấn động căn phòng khiến Fubuki giật mình choàng tỉnh. Cô bật dậy, đôi tai cáo khẽ vểnh lên đầy cảnh giác, ánh mắt mơ màng còn ngái ngủ lướt nhanh một lượt khắp văn phòng để tìm kiếm nguồn gốc âm thanh. Thế nhưng, đập vào mắt cô chỉ là một không gian hoàn toàn trống trải, tĩnh lặng như chưa từng có bất kỳ chuyện gì xảy ra.

``Lạ thật đấy\dots\ Chắc mình lại ngủ mơ rồi,'' cô lẩm bẩm, dụi dụi mắt rồi lại nằm bò ra sô-pha, ôm lấy chiếc gối ấm áp tiếp tục chìm vào giấc ngủ dang dở.

\ct

Chỉ một lát sau, tiếng bước chân rộn ràng và giọng nói líu lo vang lên từ hành lang ngoài. Cánh cửa phòng lại khẽ đẩy mở. Một cô nàng với mái tóc nâu buộc lệch tinh nghịch đang tíu tít trò chuyện cùng người tiền bối đi bên cạnh: ``Mà nè, chị Roboco-senp--''

Câu nói của cô chợt khựng lại giữa chừng. Đôi mắt to tròn mở to đầy ngỡ ngàng khi đập vào mắt cô là khung cảnh hỗn độn trước mặt. Cô nàng ngơ ngác kêu lên: ``Ể\dots?''

\chrint

Đó chính là Natsuiro Matsuri (\ruby{夏}{なつ}\ruby{色}{いろ}まつり, \ruby{Nát-xư}{Hạ}-\ruby{i-rô}{Sắc}\ Mát-xư-ri), cô nàng thuộc Thế hệ thứ Nhất, đồng nghiệp của Fubuki, kiêm thành viên mới toe của đội cổ động một trường trung học phổ thông. Mới mười sáu tuổi, Matsuri mang nguồn năng lượng tươi trẻ, hòa đồng và thân thiện đến mức có thể dễ dàng làm thân với bất kỳ ai ngay từ lần gặp đầu tiên. (Và Hikawa Kuon cũng là một trong những người may mắn lọt vào danh sách bạn bè thân thiết của cô, dẫu giữa họ có sự khác biệt hoàn toàn về giới tính).

\model{27}{r}{6cm}{matsuri}

Thế nhưng, đằng sau vẻ ngoài năng động, ngọt ngào ấy lại là một tính cách nổi loạn đầy bất ngờ. Matsuri tự nhận mình là biểu tượng của sự `trong sạch' trước công chúng, nhưng thực tế cô nàng lại sở hữu những sở thích cực kỳ nghịch ngợm sau ống kính máy quay. Cô sẵn sàng dán băng cá nhân thay vì mặc áo ngực, thích sờ soạng trêu chọc các thành viên khác khi không lên sóng, và thậm chí còn có những thói quen dở khóc dở cười, điển hình là chứng bàng quang yếu dẫn đến những tai nạn `tè dầm' khó đỡ kể từ khi cô kết thân với Inuyama Tamaki\footnote{Một nam Vtuber độc lập sở hữu diện mạo giả gái cực kỳ dễ thương và là bạn thân chí cốt của cô nàng. Thói quen khó đỡ này bắt nguồn từ thử thách chơi game Tetris cùng cậu ấy.}.

Dẫu nghịch ngợm và nông nổi là thế, Matsuri lại sở hữu một trái tim nhạy cảm, tinh tế trong việc thấu hiểu cảm xúc của người khác, khiến cô đôi khi hiện lên như một hình mẫu ``người vợ quốc dân'' chu đáo. Sự chăm chỉ, thẳng thắn khi thảo luận các vấn đề nghiêm túc cũng giúp cô ghi điểm tuyệt đối và tạo được sự kính trọng lớn từ các hậu bối trong công ty.

Khoác trên mình bộ trang phục cổ động rực rỡ, Matsuri diện chiếc áo hai dây kẻ sọc đen cá tính bên trong chiếc áo khoác lửng trễ vai màu cam nổi bật với chiếc nơ xanh lá xinh xắn. Chiếc váy xếp ly màu trắng viền cam xẻ dọc hông khoe trọn nét năng động. Đôi chân cô mang tất trắng cao cổ đi cùng giày thể thao cùng tông viền xanh lá. Đôi mắt xanh nước biển trong veo lấp lánh phản chiếu dưới mái tóc nâu tết đuôi ngựa lệch sang bên phải bằng chiếc nơ đen sọc xanh lá có đính kèm một chú thú bông nhỏ màu trắng ở giữa. Sợi ahoge ngộ nghĩnh trên đỉnh đầu cô luôn nhúc nhích theo từng cung bậc cảm xúc.

\chrint

Ánh mắt của cả hai lập tức bị hút về phía khung cửa sổ lớn. Bỏ qua sự bừa bộn thường ngày của văn phòng và dáng vẻ nằm ngủ ngon lành của Fubuki trên sô-pha, điều chấn động nhất lúc này chính là tấm kính cửa sổ lớn trước mặt đã vỡ toang, để lại một khoảng trống hun hút đón những làn gió buổi sáng luồn vào.

``Ối chà\dots\ Có ai đó đã làm vỡ tan tành cửa kính này sao?'' cô nàng có cơ thể người máy thốt lên đầy ngạc nhiên, giọng nói ngọt ngào lơ lửng giữa không trung.

\chrint
Cô nàng người máy tự xưng có ``hiệu năng tối tân'' ấy chính là Roboco-san (ロボ\ruby{子}{こ}さん, Rô-bô-\ruby{cô}{tử}-san) - thành viên đặc biệt mười bảy tuổi (theo hồ sơ đăng ký) thuộc Thế hệ thứ Không của \textit{hololive}. Đúng như cái tên của mình, cô là một thực thể người máy sở hữu công nghệ cao cấp bậc nhất, thay vì là người trần mắt thịt như các đồng nghiệp xung quanh.

\model{26}{l}{7cm}{roboco1}

Dẫu sở hữu cơ thể thép mạnh mẽ là thế, Roboco lại là định nghĩa sống của sự hậu đậu (PON), thậm chí độ vụng về của cô còn vượt mặt cả cô nàng vu nữ Sakura Miko\footnote{Nhân vật sẽ xuất hiện trở lại ở {\flaregothic ÉPISODE 4} của quyển này.}. Là một robot được lắp ghép từ những linh kiện nhặt nhạnh ở một bãi phế thải hoang vu nơi chính cô cũng chẳng nhớ nổi nằm ở đâu, việc Roboco gặp sự cố như tràn dầu động cơ\footnote{Xem trong {\flaregothic ÉPISODE 3} của quyển này.} hay vô tình làm rụng cả cánh tay\footnote{Xem trong {\flaregothic ÉPISODE 5} của quyển số Một.} đã trở thành chuyện thường ngày đối với cả văn phòng.

Tuy nhiên, đằng sau giọng nói dịu dàng, ngọt ngào như gió thoảng và vẻ ngoài ngơ ngác đáng yêu ấy lại là một cỗ máy chiến đấu vô cùng đáng sợ với bản năng sát thủ lạnh lùng. Sự ngốc nghếch, ngây thơ của cô không ít lần đẩy các thành viên khác vào những tình huống dở khóc dở cười trong suốt các tập phim.

Với bản tính thích chinh phục và chiến đấu, không ngạc nhiên khi các tựa game bắn súng là sở trường tuyệt đối của cô. Cô nàng thậm chí còn có sở thích kỳ lạ là lái xe cán qua đối thủ trong trò chơi. Tuy nhiên, điểm yếu chí mạng của nàng người máy này lại là các trò chơi kinh dị giật gân, vốn dễ dàng khiến hệ thống cảm biến của cô bị quá tải vì giật mình.

Trong những ngày đầu hoạt động, Roboco thường diện bộ trang phục vô cùng cá tính và phá cách: chiếc áo khoác biker da đen hở bạo phần ngực, thiết kế cách điệu từ quân phục với các dải sọc đỏ nổi bật trên ống tay áo. Phối cùng với đó là chiếc quần chaps kim loại bảo hộ che chắn vừa đủ, toát lên vẻ đẹp cơ khí mạnh mẽ nhưng không kém phần quyến rũ.

\chrint

\img{0-1e}

Roboco khẽ tiến lại gần khung cửa hổng, nghiêng đầu quan sát kỹ những mảnh kính còn sót lại rồi bình thản đề xuất: ``Chắc chúng ta phải đi thay ô cửa kính mới thôi nhỉ\dots\ Mà không biết họ cất mấy tấm kính dự phòng ở góc nào ta?''

Trái ngược với sự ngơ ngác của đàn chị, Matsuri chỉ biết cười trừ đầy bất lực trước đống đổ nát: ``Hỏi thế cũng hỏi, nhìn đống kính vỡ này là em biết ngay lại là tác phẩm phá hoại của Aqua-chan rồi chứ ai vào đây nữa!''

Nàng người máy nghe vậy liền gật gù tán thành: ``Có lẽ thế thật nha\~{}! Mà thôi kệ đi, chị biết chỗ cất kính dự phòng đấy! Mau nhanh tay lên nào, chúng ta còn phải sửa xong trước khi đón Kuon-senpai lên đây nữa chứ!''

``Ừ nhỉ! Hôm nay là ngày đầu tiên anh ấy đến nhận việc chính thức mà lị!'' Matsuri hào hứng hưởng ứng.

\ct

Mười lăm phút sau.

Hai người họ cùng nhau khiêng một tấm kính lớn, sáng bóng từ kho chứa về phòng. Bản thân nàng người máy sẽ dùng chính cơ thể kim loại cứng cáp của mình làm bệ đỡ để khớp tấm kính vào khung. Trong lúc vừa căn chỉnh, cả hai vừa rôm rả buôn đủ thứ chuyện trên trời dưới đất giống như những cô nữ sinh trung học:

``Roboco-senpai ơi, chừng nào em mới sở hữu được sức mạnh cơ bắp đáng nể như chị nhỉ?'' Matsuri tò mò hỏi han.

Roboco khẽ cười khúc khích, đáp lại bằng chất giọng ngọt ngào: ``Hì hì, khi nào Matsuri chịu nâng cấp cơ thể thành robot giống chị đi, bảo đảm em sẽ còn mạnh mẽ và tối tân hơn cả chị luôn đó\~{}''

Matsuri bĩu môi làm mặt xị đầy ngộ nghĩnh: ``Thôi, em xin kiếu ạ! Làm robot như chị suốt ngày bị mọi người trêu là vụng về, khéo em cũng lây bệnh PON của chị mất thôi.''

Bị trêu đúng tim đen, Roboco phồng má giận dỗi, dậm chân cãi cọ: ``Chị không có PON nhé! Rõ ràng là do hệ thống cảm biến bị nhiễu sóng thôi mà!''

Trong lúc cả hai còn đang tranh luận vô cùng rôm rả, bỗng nhiên một bóng đen nhỏ bé vo ve từ đâu bay tới, lượn lờ ngay sát mặt Roboco. Nhìn rõ sinh vật bay đáng ghét - chính là con côn trùng đã dọa Aqua nhảy lầu ban nãy - nàng Người máy lập tức giật mình hoảng hốt. Cảm biến quá tải khiến cô vung tay loạn xạ, vô tình hất tung tấm kính lớn đang cầm lên không trung. Tấm kính bay vút lên cao rồi rơi tự do xuống sàn nhà với một tiếng *RẦM* chói tai, vỡ vụn thành trăm mảnh thủy tinh lấp lánh, biến văn phòng vốn đã bừa bộn nay lại càng thê thảm hơn.

\gif{0-1f}{1}{105}

``\textbf{Áaaaa!! Con gì đáng sợ thế này!!}'' Roboco hét toáng lên đầy kinh hoàng.

``\textbf{Chờ em với, senpai ơi!!}'' Matsuri cũng hoảng hốt bám đuôi phía sau.

Chẳng thèm suy nghĩ đến hậu quả, hai cô nàng lập tức nhảy xổ ra khỏi khung cửa sổ trống hoác - tái hiện hoàn hảo cú phi thân lịch sử của Aqua trước đó - rồi nhanh chóng biến mất khỏi tầm mắt với một tốc độ nhanh đến không tưởng, để lại văn phòng chìm vào tĩnh lặng.

Ngay sau tiếng động chấn động thứ hai ấy, Fubuki lại một lần nữa giật mình choàng tỉnh khỏi giấc nồng. Cô chậm rãi ngồi dậy, đưa tay dụi dụi mắt rồi ngơ ngác nhìn quanh căn phòng trống trải. Một sự im ắng kỳ lạ bao trùm, như thể chẳng có bất kỳ ai từng đặt chân đến đây.

\img{0-1g}

Cô hướng tầm mắt về phía chiếc bàn đối diện. ``Ủa? Cái gì thế kia?''

Fubuki tò mò rảo bước tiến lại gần, phát hiện ra một chiếc máy quay nhỏ đang đặt âm thầm trên bàn, ống kính vẫn đang ghi hình mọi cử chỉ từ lúc cô bước vào phòng đến giờ. ``Sao lại có máy quay đặt ở đây nhỉ\dots?''

Tò mò, cô khẽ nâng tay xoay nhẹ ống kính máy quay hướng lên trần nhà. Ngay lập tức, toàn bộ bề mặt trần phòng hiện ra trước ống kính, phủ kín bởi những chữ ký quen thuộc của các thành viên \hll\ cùng những nét vẽ nguệch ngoạc vô cùng đáng yêu. Thế nhưng, đập vào mắt người xem chính là dòng chữ lớn được viết nổi bật đầy hào hứng ở chính giữa: ``Coming Soon!!''

Nàng Cáo trắng nghiêng đầu suy nghĩ một lát, đôi tai cáo khẽ vẩy nhẹ: ``Sắp tới có dự án gì đặc biệt sao?'' Ánh mắt cô ánh lên nét tò mò khó giấu, cô nhún vai cười mỉm: ``Chẳng lẽ mọi người đang âm thầm chuẩn bị một bất ngờ lớn mà mình chưa biết nhỉ? Nhưng thôi kệ đi, chắc chắn sẽ là một dự án cực kỳ thú vị đây!''

\img{0-1h}

%============================================TRUYỆN PHỤ=======================================
\omk

\pov{Hikawa Kuon}

``Chỗ này\dots\ đúng không nhỉ?''

Tôi đứng chững lại giữa ngã tư đường, ngơ ngác nhìn quanh lối ra của ga tàu điện ngầm Shinjuku đông đúc tại Tokyo. Cầm trên tay tấm bản đồ ga tàu chằng chịt như mạng nhện, tôi chỉ biết thở dài bất lực. Ở Kawachi, tôi vốn chỉ quen đi xe buýt, cứ nhảy lên xe là xong, chứ còn hệ thống tàu điện ngầm ở cái thành phố khổng lồ này quả thực là một cơn ác mộng đối với một kẻ mù đường như tôi. Tính cả thời gian đi tàu điện ngầm và cuốc bộ rạc cả cẳng, tôi đã mất gần ba tiếng đồng hồ chỉ để tìm đường. Địa chỉ trên tờ hướng dẫn viết loằng ngoằng quá thể! Mà không biết đây có thực sự là nơi làm việc sắp tới của mình không nữa.

Sau buổi sáng chuẩn bị kỹ lưỡng, hôm nay tôi sẽ đến nơi này lần đầu tiên để chính thức nhận công việc. Tôi tự nhủ bản thân phải nắm rõ mọi ngóc ngách để không bị bỡ ngỡ khi bước vào môi trường mới vào ngày mai. Để chuẩn bị cho ngày trọng đại này, tôi đã cố gắng ăn diện sao cho thật lịch sự - tất nhiên là theo tiêu chuẩn khắt khe của bố tôi, chứ tôi thì vốn xuề xòa. Đó là một chiếc áo thun đỏ đơn giản phối cùng quần jeans xanh sẫm cứng cáp. Lúc thắt chiếc dây lưng da, tôi phải hơi hóp bụng lại một chút để chốt khóa khớp vào lỗ. Bố đứng bên cạnh nhìn thấy, liền tặc lưỡi chê bai: ``Đấy, nhìn cái bụng mỡ chưa kìa! Mới hai mươi mấy tuổi đầu mà bụng đã phệ ra như ông chú trung niên rồi. Lo mà đi tập thể dục đi con ạ, suốt ngày nằm ườn ra đấy lười chảy thây!''

Tất nhiên, lời khuyên ấy như gió thoảng qua tai, tôi chỉ ậm ừ cho qua chuyện vì vốn lười vận động. Dù bố có nhắc nhở cả chục lần đi chăng nữa, tôi vẫn thà dành thời gian gõ code còn hơn chạy bộ. Để giữ đúng phong cách thoải mái, tôi từ chối đi giày vì sợ bí chân, thay vào đó là chọn đôi dép quai hậu da màu đen thông thoáng. Mái tóc đen cũng được chải chuốt gọn gàng trước gương, dẫu vẫn còn vài sợi tóc xoăn cứng đầu dựng ngược lên chống đối.

Sau một hồi cuốc bộ dọc theo con phố nhỏ, tôi chợt nhìn thấy một tòa nhà cao khoảng chín, mười tầng hiện lên phía trước. So với những tòa nhà văn phòng cao chọc trời xung quanh, nó trông khá khiêm tốn, nhưng màu sơn xanh lam nổi bật cùng dòng chữ logo \textit{hololive} to tướng ở mặt tiền đã xác nhận tất cả. Chắc chắn không thể nhầm lẫn đi đâu được nữa.

\ct

Tôi hít một hơi thật sâu rồi mò mẫm bước vào sảnh tòa nhà, tiến thẳng về phía văn phòng nơi các thành viên \hll\ làm việc. Thế nhưng, vừa đẩy cửa bước vào trong, một cảnh tượng dở khóc dở cười lập tức đập vào mắt tôi. Ba cô gái đang ôm chầm lấy nhau thành một cục, run lẩy bẩy như cầy sấy giữa căn phòng văn phòng hỗn độn.

Tôi ngơ ngác gãi đầu, cất tiếng hỏi đầy thắc mắc: ``Cái gì thế này?''

Vừa nghe thấy giọng tôi, cả ba người họ như bắt gặp được vị cứu tinh, đồng loạt buông nhau ra rồi cuống cuồng chạy ào tới. Họ chính là Minato Aqua, Natsuiro Matsuri và Roboco, ba cô gái mà tôi mới chỉ trao đổi công việc sơ qua qua cuộc gọi Zoom tối qua mà chưa có cơ hội gặp mặt ngoài đời.

``Ơn giời! Kuon-senpai, anh đây rồi\dots!'' Aqua mếu máo kêu lên, gương mặt tái mét.

Matsuri cũng luống cuống tiếp lời, hai tay khua loạn xạ: ``Kuon-senpai, cứu bọn em với!''

``Đúng đó, bọn em vừa thấy một con quái vật bay cực kỳ đáng sợ lao vào phòng á!'' Roboco hớt hải bổ sung, đôi mắt người máy chớp liên hồi vì hoảng loạn.

Tôi đứng đực ra đó, đầu óc mơ hồ chưa kịp hiểu mô tê gì. ``Quái vật bay? Mấy đứa đang nói cái gì thế?''

Nhưng họ chẳng để tôi kịp định thần, cứ thế tranh nhau giải thích. ``Nó to lắm luôn á, trông cứ như một con ong khổng lồ vậy!'' Aqua giơ hai tay mô tả kích thước đầy phóng đại.

``Bọn em sợ nó đốt lắm, anh làm gì đó diệt nó đi chứ!'' Matsuri hét lên đầy khẩn thiết.

Chưa kịp để tôi mở miệng trấn an, một tiếng vo ve chói tai đột ngột vang lên sát sườn. Sinh vật bay bí ẩn - thủ phạm gây ra vụ hỗn loạn ban nãy - bất ngờ lượn từ góc phòng ra và đậu ngay trên bức tường phía sau lưng tôi.

Cả ba cô gái lập tức rú lên, đồng thanh hét lớn cảnh báo: ``\textbf{Đằng sau anh kìa!}''

Tôi bình thản quay đầu lại nhìn. Hóa ra chỉ là\dots\ một con gián đất to xác. Nhanh như cắt, tôi rút ngay chiếc dép lào mang theo bên mình - chao ôi, chẳng hiểu sao tôi lại có linh tính mang theo món vũ khí huyền thoại này nữa - rồi vung tay làm một cú vụt cực mạnh.

*Bép!*

Một tiếng động dứt khoát vang lên. Chỉ trong chớp mắt, sinh vật đáng sợ kia đã còng queo dưới sàn nhà, chính thức chầu trời không kịp ngáp.

``Xời. Mấy con này ở nhà anh đầy, quét đi là xong. Dễ như ăn kẹo,'' tôi phủi phủi tay, đút dép lại vào trong túi đầy tự đắc.

Aqua rụt rè tiến lại gần, mắt ti hí nhìn xuống sàn: ``N-Nó chết thật chưa anh?''

Nhìn thấy thứ sinh vật kia không còn cựa quậy nữa, họ mới đồng loạt thở phào nhẹ nhõm.

``Phù\dots\ Được cứu rồi. May mà có anh ở đây,'' Matsuri vuốt ngực tự an ủi.

Roboco cũng gật gù đầy cảm kích: ``Cảm ơn anh nhiều nhé, Kuon-senpai!''

Tôi chưng hửng nhìn cô nàng người máy, không nhịn được mà lên tiếng mỉa mai: ``Mà này Roboco, em là robot cơ khí cơ mà, bị ong đốt thì thấm tháp gì mà phải sợ? Cái tia laser tích hợp trên tay của em để làm cảnh à?''

Bị khích tướng, Roboco lập tức phồng má phản ứng lại đầy giận dỗi: ``Anh điên à! Bắn laser ở đây thì có mà cháy cả cái văn phòng này! Anh muốn bọn em phải đền tiền cả phòng hả?''

``Ờ thì, cứ diệt được nó là được chứ gì. Đền bao nhiêu anh trả,'' tôi bật cười đáp lại bằng một câu đùa cợt. Tất nhiên, với độ PON trứ danh của mình, Roboco chẳng thể nhận ra đó là lời nói đùa, cô nàng lập tức xua tay lia lịa: ``Thôi, em xin kiếu!'' Matsuri đứng bên cạnh cũng gật đầu lia lịa đồng tình.

Nói đùa thế thôi chứ tôi đào đâu ra tiền mà đền. Nếu ví tiền rực rỡ thì tôi đã chẳng phải lặn lội sang đây xin việc làm bán thời gian đúng chuyên ngành để kiếm mức lương cao ngất ngưởng kia.

Sau khi giải quyết xong mối họa côn trùng, bốn chúng tôi cùng nhau bước vào văn phòng chính. Tại đây, Fubuki vẫn đang loay hoay săm soi chiếc máy quay nhỏ đặt trên bàn.

``A, mọi người đến rồi! Chào buổi sáng cả nhà!'' Nàng Cáo trắng vui vẻ vẫy tay chào khi thấy chúng tôi bước vào.

Cả ba cô gái đi cùng tôi liền đồng thanh đáp lại lời chào.

Tôi tiến về phía bàn, chợt nhận ra điều gì đó lạ lùng: ``À, cái máy quay này trông quen thế nhỉ? Hình như là của chị A-san đấy. Chị ấy lắp ở đây làm gì ta\dots?'' Tôi cầm chiếc máy quay lên kiểm tra, phát hiện ra màn hình hiển thị của nó đã tắt ngúm tự bao giờ.

``Lạ thật đấy, tự dưng ai lại đặt máy quay ở đây làm gì chứ?'' Matsuri thắc mắc gãi đầu.

Aqua cũng nghiêng đầu đăm chiêu: ``Chắc là A-chan đặt nó ở đây để quay thử cảnh gì đó chăng? Nhưng sao tự dưng nó lại tắt rồi nhỉ?''

Roboco nhún vai: ``Chị cũng không rõ nữa. Nhưng A-chan hay bày ra mấy trò bất ngờ lắm, biết đâu đây là một phần trong kế hoạch bí mật nào đó thì sao.''

``Hoặc đơn giản là máy hết pin thôi. Anh đoán thế,'' tôi phỏng đoán, đặt chiếc máy quay trở lại bàn rồi thở dài: ``Nhưng vấn đề lớn nhất ở đây là, tại sao chị ấy lại làm thế?''

Fubuki gật đầu phụ họa đầy hào hứng: ``Đúng vậy á! Tại sao chị ấy lại lén lút đặt nó ở đây mà không nói với chúng ta một lời nào hết nhỉ? Cứ như là\dots\ muốn tụi mình tự đi tìm câu trả lời vậy!''

``Cái trò này nghe có vẻ giống một dạng thử thách hay trò chơi mật mã ghê nha,'' Matsuri tủm tỉm cười.

Aqua quay sang nhìn tôi, đôi mắt lấp lánh sự kỳ vọng: ``Anh nghĩ sao về chiếc máy này hả Kuon-senpai?''

Tôi xua tay, quyết định giữ thái độ trung lập: ``Anh nghĩ tốt nhất chúng ta không nên táy máy vào nó làm gì. Cứ để chiếc máy đó yên vị ở đấy, đợi chị A-san tới giải thích sau là rõ ngay thôi.''

Nói rồi, tôi cất chiếc máy quay gọn vào góc bàn, lập tức chuyển chủ đề hỏi bốn người về dự án tiếp theo của họ: ``Mà này, dự án cái gì mà\dots\ \textit{Holo no Graffiti} ấy, tiến hành đến đâu rồi? Sáng nay anh nghe chú Tanigo loáng thoáng nhắc tới tên này qua điện thoại\dots''

Cả bốn cô gái đồng loạt nhìn nhau, rồi quay ngoắt sang nhìn tôi với vẻ mặt ngơ ngác đến tột độ.

``Holo no Graffiti? Em chưa từng nghe cái tên này bao giờ luôn á\dots'' Matsuri chớp chớp mắt.

Aqua cũng tò mò hỏi ngược lại: ``Là cái gì thế Kuon-senpai? Một dự án mới của Cover Corp hả anh?''

Roboco gãi má đăm chiêu: ``Chắc chắn không phải là chương trình nào mà chúng ta từng tham gia trước đây rồi đúng không?''

Fubuki cũng vểnh tai cáo lên, cố lục lọi ký ức: ``Chúng mình có biết cái tên này không ta? Em cũng không nhớ là đã từng nghe qua nữa.''

Nhìn phản ứng ngơ ngác của họ, tôi không khỏi ngạc nhiên. Hóa ra cái tên \textit{Holo no Graffiti} lại là một thứ gì đó hoàn toàn xa lạ đối với chính các Vtuber trong công ty. Thực ra, bản thân tôi cũng chẳng hiểu rõ ngọn ngành về nó, chỉ nghe qua loa từ cuộc điện thoại chớp nhoáng sáng nay với chị A-san và chú Tanigo. Có lẽ do họ giải thích quá vắn tắt, hoặc do lúc ấy đầu óc tôi đang mải nghĩ về chuyện khác nên không chú ý.

Tôi gãi đầu cười trừ, đành phải thú nhận: ``Thú thật là anh cũng chẳng rõ đầu đuôi dự án này ra sao nữa\dots\ A-san và chú Tanigo chỉ nói sơ qua trong cuộc gọi sáng nay thôi. Hình như nó là một chuỗi video hoạt hình ngắn hay cái gì đó liên quan đến nghệ thuật số, nhưng chi tiết cụ thể thì anh chịu.''

Matsuri nghe vậy liền lay lay tay tôi, nũng nịu trêu chọc: ``Ối chà\~{}, thế thì senpai phải đi tìm hiểu rồi giải thích cho tụi em nghe đi chứ!''

``Đúng đó! Có thông tin gì mới anh phải chia sẻ ngay cho mọi người biết với nha!'' Aqua cũng hùa theo.

Roboco hào hứng gật đầu: ``Em muốn biết thêm nữa, nghe cái tên thôi đã thấy giống một dự án hoành tráng rồi!''

Tôi chỉ biết dở khóc dở cười trước sự nhiệt tình của ba cô nàng: ``Nhưng mà anh đã bảo là anh có biết cái quái gì nữa đâu cơ chứ\dots''

Nhìn vẻ mặt bối rối của tôi, bốn cô gái chỉ biết cười xòa. Tôi cố gắng vớt vát lại chút uy nghiêm của người quản lý: ``Để anh thử hỏi lại xem sao\dots\ Nhưng hiện tại anh chỉ biết tên dự án và việc nó sắp sửa được triển khai thôi. Nghe bảo là sẽ bắt đầu chạy chính thức từ tuần sau, còn hôm nay chỉ là buổi ghi hình demo chạy thử nghiệm thôi.''

``Vậy à\dots'' Họ nhìn nhau với ánh mắt lấp lánh sự tò mò và phấn khích: ``Vậy thì chúng ta bắt tay vào triển khai thôi chứ nhỉ?''

Cả bốn cô gái đồng thanh hô vang đầy khí thế: ``Vâng ạ\~{}!''

Tôi mỉm cười gật đầu, lòng tràn đầy quyết tâm: ``Tốt lắm. Vậy thì\dots\ chúng ta khởi hành thôi! Let's roll!''

Và như thế, cuộc phiêu lưu của tôi - Hikawa Kuon - cùng với các thành viên tinh nghịch của gia đình \textit{hololive} đã chính thức bắt đầu từ ngày hôm đó.

\dots\ Ít nhất là sau khi tôi hoàn tất thủ tục nhận việc chính thức từ chú CEO Tanigo đã.

\ct

Dưới sự dẫn đường của chị Yuujin A, tôi tìm đến phòng làm việc của CEO nằm ở tầng trên của tòa nhà. Tại đây, tôi cùng chú Tanigo và chị ấy ngồi lại với nhau để thống nhất về lịch trình làm việc cụ thể tại công ty.

Theo lịch trình đã thống nhất, tôi sẽ làm việc bán thời gian từ thứ Hai đến thứ Sáu. Ca sáng bắt đầu từ tám giờ đến mười một giờ hai mươi lăm phút, còn ca chiều từ một giờ đến sáu giờ tối. Quãng đường đi từ nhà đến công ty mất khoảng bốn mươi phút đi xe buýt và thêm mười phút đi bộ, nên trừ những ngày vướng lịch học ở trường, tôi luôn tự giác có mặt sớm trước cả tiếng đồng hồ để chuẩn bị.

Riêng hai ngày cuối tuần, tôi đăng ký làm toàn thời gian từ tám giờ ba mươi sáng đến năm giờ chiều. Nhưng tất nhiên, tôi sẽ cố gắng nán lại đến sáu giờ tối mới chịu ra về. Các bạn có thể gọi tôi là kẻ cuồng công việc cũng chẳng sai, nhưng thực ra tôi chỉ muốn mau chóng thích nghi với môi trường mới mà thôi.

Trong buổi trò chuyện, chú Tanigo - hay còn gọi là YAGOO - cùng chị A cũng hỏi thêm một số thông tin để hiểu rõ hơn về tôi. Tôi thẳng thắn chia sẻ rằng điểm mạnh của mình là sự trung thực, đôi khi thẳng tính đến mức hơi thô thiển, và thái độ làm việc hết sức cần cù. Để chứng minh cho trí thông minh logic mà bản thân tự nhận là khá tốt, tôi đã chủ động viết thử một vài đoạn code mẫu để hai người cùng xem qua. Tôi cũng cho biết mình có khả năng giao tiếp bằng tiếng Anh. Thực chất, tôi còn nói được nhiều thứ tiếng khác nữa, nhưng tôi quyết định giấu nhẹm đi. Dù sao thì ở một công ty giải trí Nhật Bản thế này, chắc cũng chẳng mấy khi dùng đến\dots\ đúng không nhỉ?

Sau khi các thủ tục giấy tờ được ký kết hoàn tất, mọi thứ coi như đã an bài. Kể từ giây phút này, tôi chính thức trở thành tổng quản lý của văn phòng \hll\ - một công việc mà lúc bấy giờ tôi vẫn chưa hình dung hết khối lượng và tính chất của nó, nhưng chắc chắn là vị trí trong mơ của biết bao gã đàn ông ngoài kia. Đồng hành cùng tôi trên chặng đường sắp tới sẽ là người đồng nghiệp đáng mến, chị Yuujin A.

Còn về dự án \textit{Holo no Graffiti} đầy bí ẩn kia thì sao? Tôi đã thử gặng hỏi cả vị CEO đáng kính lẫn người đồng nghiệp mới của mình, nhưng họ chỉ nhìn tôi mỉm cười đầy ẩn ý rồi đồng thanh đáp: ``Cậu hãy tự mình trải nghiệm đi rồi sẽ rõ.''

Giờ thì tôi chỉ biết tự cầu chúc cho bản thân chân cứng đá mềm để xứng đáng với trọng trách mới này, và quan trọng nhất là\dots\ kiếm thật nhiều tiền mang về phụ giúp bố mẹ\dots? Ừm, hy vọng là thế.

Mà thật lòng mà nói\dots\ tôi vừa tự đẩy mình vào cái tình huống quái đản gì thế này cơ chứ?

\end{document}